\section{Methods}
\label{sec:methods}

\subsection{Construction of the bilinear forms}
\label{app:construction}

\paragraph{Polynomial models.} Represent a polynomial vector field on $\R^n$ by
an $n \times m$ array of coefficients over the monomials of total degree at
most $d$. Multiplication and differentiation of monomials are index tables
computed once, so the bracket~\eqref{eq:bracket} of two such fields is exact
integer and rational linear algebra. Fixing the candidate basis
$\zeta_1,\dots,\zeta_P$ and the model basis $\{x^a e_i\}$ and computing the
bracket of every pair gives a three-index tensor $T \in \R^{J \times P \times
Q}$, with $J$ the number of coefficients of a bracket, such that
$\mathrm{coef}([\xi_\theta,F_\beta]) = T \cdot \theta \cdot \beta$, and likewise
$T^{(1)}, T^{(2)}$ for the two individual terms, with $T = T^{(1)}-T^{(2)}$.
Multiplying the first index by $W^{1/2} \otimes g_\nu^{1/2}$, where $g_\nu$ is
the monomial Gram matrix of $\nu$ and $W$ the output metric, converts Euclidean
norms of contracted vectors into the norms $\|\cdot\|_{\nu,W}$ of the
corresponding fields. Then $L_\theta = T_w \cdot \theta$, and the columns of
$A(\beta) = T_w \cdot \beta$ give $C(\beta) = A(\beta)^{\!\top}A(\beta)$.
The matrix $M(\beta)$, defined by $M(\beta)\theta = M_\theta\beta$, stacks
$T^{(1)}_w\cdot\beta$ and $T^{(2)}_w\cdot\beta$ and gives $D(\beta) =
M(\beta)^{\!\top}M(\beta)$. The matrix
$g_\nu$ is available in closed form for the uniform measure on a box (with
equal or unequal half-widths), the Gaussian, and the uniform measure on a ball
or annulus (the last two by separating radial and angular moments), so no
quadrature error enters. For the \emph{data} measure we substitute the
empirical Gram matrix of the same monomial basis; the difference between the
two is exactly what \Cref{prop:kappa} measures. To evaluate the extrapolation
factor, $G_\mu$ is symmetrised and treated as singular when its smallest
eigenvalue is at most $10^{-10}$ times its largest, in which case
$\sqrt{\kappa\kappa'}$ is reported as infinite; otherwise $\kappa$ and
$\kappa'$ are computed with symmetric inverse square roots. No pseudoinverse or
ridge is used.

\paragraph{Dictionary models.} For $F_\beta = \sum_{i,a}\beta_{i,a}\phi_a e_i$
with differentiable $\phi_a$, the two terms of the bracket of $\xi$ with
$\phi_a e_i$ are $(\nabla\phi_a\cdot\xi)\,e_i$ and $\phi_a\,(D\xi)\,e_i$, both
linear in $\beta$ pointwise. Evaluating them at the nodes $x_k$ of a quadrature
rule with weights $w_k$ for $\nu$ and stacking $\sqrt{w_k}\,W^{1/2}$ times
their values gives tensors with the same role as $T^{(1)}_w, T^{(2)}_w$; their
output index is compressed by an orthonormal basis of the span of all
contracted vectors, which leaves every norm unchanged. A uniform grid of $p$
points per angle integrates trigonometric polynomials of degree below $p$
exactly, and a Gauss--Legendre rule with $p$ nodes integrates polynomials of
degree up to $2p-1$ exactly. For the radial library we integrate in polar
coordinates, Gauss--Legendre (40 nodes) in the radius with the area weight and
a uniform grid (24 points) in the angle, and for the pendulum library a
$32 \times 32$ Gauss--Legendre product rule. For the radial family and the
pendulum the test suite checks that doubling both rules changes no defect by
more than a relative $10^{-10}$, and that Monte Carlo integration of the
analytic bracket over the target measure ($10^{6}$ points) agrees with the
quadrature defects to within three standard errors. With a monomial dictionary and an exact rule the construction
reproduces the closed-form tensors to $10^{-10}$, which the test suite also
checks.

\subsection{Certified numerics}
\label{sec:numerics}

All extrema are over the ball $\|u\| \le R$ that parameterises the ellipsoid.
For $\max\|a + Bu\|$ we use the dual bound $\|a+Bu\|^2 \le a^{\!\top}a + \omega R^2
+ c^{\!\top}(\omega I - B^{\!\top}B)^{-1}c$, valid for every multiplier $\omega >
\lambda_{\max}(B^{\!\top}B)$ with $c = B^{\!\top}a$; for $\min\|a+Bu\|$ the bound
$\|a+Bu\|^2 \ge a^{\!\top}a - c^{\!\top}(B^{\!\top}B + \omega I)^{-1}c - \omega R^2$,
valid for every $\omega \ge 0$ with the matrix positive definite (and the
least-squares residual when the unconstrained minimiser lies in the ball). For
the indefinite problems of \Cref{prop:closedform}(ii) we use the dual function
$\psi(\omega) = s_0 - g^{\!\top}(H + \omega I)^{-1}g - \omega R^2$ of
$\min\,u^{\!\top}Hu + 2g^{\!\top}u + s_0$, valid and concave for every
$\omega \ge 0$ with $H + \omega I$ positive definite (at the end point $\omega =
-\lambda_{\min}(H) > 0$ it remains valid when $g$ is orthogonal to the bottom
eigenspace of $H$; otherwise that end point is approached as a limit). In each case the multiplier is located by a safeguarded Newton iteration
on the derivative of the dual (the secular equation), in the reciprocal form of
Mor\'e and Sorensen, which is concave and nearly linear, with bisection as a
fall-back \citep{gander1989constrained,more1983computing,conn2000trust};
whatever multiplier is returned, the value is a valid bound, and the minimiser
of the subproblem (with a component along the bottom eigenvector in the hard
case) is a feasible point. The sharp certificate uses this in Dinkelbach's
iteration for fractional programmes \citep{dinkelbach1967nonlinear}: starting
from the plug-in defect, each step solves the indefinite subproblem at the
current level $t$ and moves $t$ to the squared ratio at its minimiser, a
feasible point, so the iterates increase to $U(\theta)^2$ from inside; the
square root of the final value, inflated by a relative $10^{-7}$, is accepted
only when the dual
bound certifies $\min_u\, t\|M_\theta\beta(u)\|^2 - \|L_\theta\beta(u)\|^2 > 0$
there, and bisection between the attained value and the separate bound takes
over otherwise. A further relative margin of $10^{-9}$ covers rounding, and the
analogous construction from above gives $L(\theta)$. A denominator whose lower
bound is below $10^{-12}$ of its scale is treated as vanishing, and the
generator as not certifiable. The separate bound uses
$\sup_{\mathcal B}\|L_\theta\beta\|$ and $\inf_{\mathcal B}\|M_\theta\beta\|$ (and,
when model error is admitted, the two denominator blocks separately, because
the model-error terms perturb them by different amounts). One sharp
certificate costs about $3$~ms of CPU time for $Q = 20$ and $10$~ms for
$Q = 60$ on one core.

For subspaces, with $V$ also denoting a matrix whose orthonormal columns span
the subspace and $\theta$ ranging over its unit vectors, we bound
$\sup_{\theta}\|A(\beta)\theta\| \le
\sigma_{\max}(A(\bhat)V) + \sup_{\mathcal B}\|(A(\beta)-A(\bhat))V\|_F$ and
$\inf_{\theta}\|M(\beta)\theta\| \ge \sigma_{\min}(M(\bhat)V) -
\sup_{\mathcal B}\|(M(\beta)-M(\bhat))V\|_F$, each supremum being the norm of a
linear image of $\beta$ and hence one more trust-region maximum. The refutation
bound uses the same construction with $V = I$, the whole candidate class:
$\min_\theta\defect_\nu \ge (\sigma_{\min}(A(\bhat)) - r_A)_+ /
(\sigma_{\max}(M(\bhat)) + r_M)$ on $\{\beta^\star \in \mathcal B\}$, with $r_A =
\sup_{\mathcal B}\|A(\beta)-A(\bhat)\|_F$ and $r_M = \sup_{\mathcal
B}\|M(\beta)-M(\bhat)\|_F$. Generalised eigenproblems are solved on the range
of $D(\bhat)$; directions in its null space (relative eigenvalue at most
$10^{-12}$) have an undefined defect and are never declared or certified.

\paragraph{Pointwise bound.} For a generator fixed in advance, or chosen on
one half of the data and bounded with the fit to the other half, the error
probability is split into $0.98\alpha$ for the numerator, $0.01\alpha$ for the
denominator and $0.01\alpha$ for the noise scale: $\hat\Sigma$ is multiplied by
the upper confidence factor $\mathrm{dof}/\chi^2_{\mathrm{dof};\,0.01\alpha}$
(1 for the asymptotic regions), and $\|L_\theta\bhat\|$ and $\|M_\theta\bhat\|$ are each moved up and down by
the Gaussian-concentration width $(\mathrm{tr}\,\Gamma)^{1/2} +
(2\|\Gamma\|_{\mathrm{op}}\log(1/\alpha'))^{1/2}$ at their share $\alpha'$,
with $\Gamma$ the inflated covariance of $L_\theta\bhat$ or $M_\theta\bhat$.
The bounds on the defect are the ratios of these bounds. When model error is
admitted, the denominator's share is spent on its two blocks, $0.005\alpha$
each.

\paragraph{Closure diagnostic.} For a certified plug-in subspace $V_r$ with
$r \ge 2$ we take a basis $\xi_1,\dots,\xi_r$ that is orthonormal in the
$\|\cdot\|_{\nu,W}$ geometry of generator fields and compute each commutator
$[\xi_a,\xi_b] = D\xi_b\,\xi_a - D\xi_a\,\xi_b$ exactly. For affine generators
it lies in the candidate class, which is closed under the bracket. The
reported closure defect is the largest over pairs of
\[
  c(a,b) = \frac{\|P_\perp[\xi_a,\xi_b]\|_{\nu,W}}
  {\bigl(\|D\xi_b\,\xi_a\|^2_{\nu,W} + \|D\xi_a\,\xi_b\|^2_{\nu,W}\bigr)^{1/2}},
\]
with $P_\perp$ the orthogonal projection, in the same geometry, onto the
complement of $V_r$ in the candidate class. For a class closed under the
bracket, as here, $c = 0$ for every pair exactly when $V_r$ is a Lie
subalgebra, and $c \le \sqrt 2$. A commutator whose norm
exceeds $5\%$ of the same denominator is itself a candidate generator, and its
defect is bounded by the sharp certificate from the same region at the same
level; the commutators count as certified when every such bound is at most
$\delta$. Smaller commutators, as in an abelian algebra, are reported as
negligible. For a candidate class that is not closed under the bracket
(quadratic generators, for example), a commutator need not be a candidate
generator and then cannot be certified, and $c$ measures only its part inside
the class: the part outside it must be reported alongside $c$, and $c = 0$
alone does not establish closure.

\subsection{Confidence regions}
\label{sec:regions}

Unless stated otherwise the region is the exact least-squares ellipsoid of
\Cref{prop:closedform}(iii) (``OLS'' below). The sandwich regions use
$\hat\Sigma = (I_n\otimes (\Phi^{\!\top}\Phi)^{-1})\,\Omega\,(I_n \otimes
(\Phi^{\!\top}\Phi)^{-1})$ with scores $z_t = r_t \otimes \phi(x_t)$, where $r_t \in \R^n$ is the
least-squares residual of row $t$. White's
estimator (HC0) takes $\Omega = \sum_t z_t z_t^{\!\top}$ and the radius $R^2 =
\chi^2_{Q;\,1-\alpha}$. The cluster-robust estimator is the jackknife (CR3) form
recommended when clusters are few \citep{cameron2015practitioner}: with $n_{\mathrm c}$
trajectories, $\Omega = \frac{n_{\mathrm c}-1}{n_{\mathrm c}}\sum_g s_g s_g^{\!\top}$, where $s_g =
\Phi_g^{\!\top}(I - H_{gg})^{-1} r_g$ collects the residuals of trajectory $g$
inflated by its leverage block $H_{gg} =
\Phi_g(\Phi^{\!\top}\Phi)^{-1}\Phi_g^{\!\top}$; it uses the small-sample radius
$R^2 = Q\,\frac{n_{\mathrm c}-1}{n_{\mathrm c}-Q}\,\mathrm{F}_{Q,\,n_{\mathrm c}-Q;\,1-\alpha}$,
the Hotelling form, which tends to the chi-square quantile as $n_{\mathrm c}$
grows. The Newey--West estimator adds, within each trajectory, the products of
scores at lags $j = 1,\dots,\ell$ with Bartlett weights $1 - j/(\ell+1)$, with
$\ell$ chosen by Andrews' AR(1) plug-in rule from the pooled lag-one
autocorrelation of the residuals, and uses the chi-square radius
\citep{newey1987simple,andrews1991heteroskedasticity}. Directions outside the
range of $\hat\Sigma$ (relative eigenvalue at most $10^{-12}$), such as those
missed by a cluster-robust estimate with fewer trajectories than coefficients,
and the coefficient directions that a rank-deficient design cannot see are
treated as unidentified, and any certificate that depends on them is
infinite. Region coverage is checked by evaluating
$(\bhat-\beta^\star)^{\!\top}\hat\Sigma^{-1}(\bhat-\beta^\star) \le R^2$ against
the known $\beta^\star$. The estimators used are: OLS in every experiment
except the following; White and OLS for the error-distribution study; White and
cluster-robust for the trajectory study; OLS, White, cluster-robust and
Newey--West for the serial-dependence study and for the drifter case study;
White for the lynx--hare case study, where with $19$ rows it carries no
asymptotic justification.

\subsection{Model-class procedures}
\label{sec:classproc}

The lack-of-fit test is the $F$-test of the smaller class against the larger
on the same rows, with residual sums of squares from both fits,
$n(m_{\mathcal H}-m)$ and $n(N - m_{\mathcal H})$ degrees of freedom, and
level $0.05$. Its denominator is the residual mean square of the largest
class, and the certificate uses the residual variance of the class in which it
is computed, so the two share one estimate when the gate selects the largest
class and use different estimates from the same rows otherwise; under
\emph{Split} they use disjoint halves of the rows, and \emph{Largest} and
\emph{Simultaneous} run no test. The
\emph{gate} chooses the smallest class whose test against the largest does not
reject and certifies there; no error budget is allocated to the test, because
the gate carries no guarantee. \emph{Split} applies the gate to the first half
of the rows and certifies on the second. \emph{Largest} certifies in the
largest class at level $\alpha$. \emph{Simultaneous} implements
\Cref{prop:classes} with $\alpha/K$ per class; the omitted part of class $k$
is computed from the least-squares projection of the largest class on class
$k$ over the design, and its $\|\cdot\|_{\nu,W}$ norm is maximised over the
largest-class region. The invisible defect of \Cref{prop:lof} is evaluated
with $\mathcal H$ the largest class, $\mu_N$ replaced by $\nu$ (the designs are
uniform on the target box), $\gamma = 0.5$ and the denominator
of~\eqref{eq:invisible} evaluated at $F_{\beta^\circ}$ (the degree-3 Hopf
field), that is, to first order in $h$. The resulting values describe one
omitted part, $h^\star$, and are not upper bounds; evaluating the denominator
exactly at $F_{\beta^\circ} \pm h^\star$ gives $0.039$--$0.040$ (degree 5) and
$0.085$--$0.089$ (degree 7) at $N = 3200$, depending on the sign of $h^\star$.

\subsection{Baselines}
\label{app:baselines}

\emph{Fixed threshold}: declare the span of all generalised eigenvectors of
$(C(\bhat), D(\bhat))$ whose plug-in defect is at most $\tau$; sweeping $\tau$
traces the whole family, including $\tau = \delta$. \emph{Eigengap}: declare
the $r$ maximising $\defect_{(r+1)}/\defect_{(r)}$ over $r \ge 1$; this rule
cannot express ``no symmetry'', which is one reason its error rate is high.
\emph{Weyl perturbation certificate}: bound $\lambda_k(C(\beta^\star)) \le
\lambda_k(C(\bhat)) + \|C(\beta^\star)-C(\bhat)\|$ with the perturbation
bounded by $2\|A(\bhat)\|\,\|\Delta A\|_F + \|\Delta A\|_F^2$, $\|\Delta A\|_F$
maximised over the same region, and divide by a lower bound on
$\lambda_{\min}(D)$; this is the natural singular-value perturbation answer
\citep{bhatia1997matrix,davis1970rotation,yu2015useful}, valid
and loose because it spends one global perturbation budget on every direction.
\emph{Rank test}: for each $r$, test $H_0$: the $r$ smallest plug-in directions
span an exact nullspace, using the squared norm of the restricted bracket
referred to a Welch--Satterthwaite approximation of its null distribution, and
report the largest $r$ not rejected, a Wald-type rank test in the spirit of
Robin and Smith~\citep{robin2000tests} and Kleibergen and Paap~\citep{kleibergen2006generalized}. \emph{Split
significance test}: select the smallest-defect direction on one half of the
data, test $H_0$: exact symmetry on the other half by the same statistic, and
declare a symmetry when the test fails to reject. \emph{\symcert{} with sample
splitting}: fit the first and the second half of the rows separately, take the
plug-in generalised eigenvectors of the first-half fit in order of plug-in
defect, certify the span of the first $r = 1, 2, \dots$ of them with the
simultaneous certificate computed from the second-half region (the sharp
certificate for $r = 1$, the subspace bound otherwise), and declare the
largest $r$ before the first span whose bound exceeds $\delta$. In the
dimension study only the first direction is bounded, with the pointwise bound.
\emph{Parametric bootstrap}: draw $\beta^* \sim N(\bhat,\hat\Sigma)$ $400$
times and report the $\alpha$ and $1-\alpha$ quantiles of
$\defect_\nu(\xi_\theta;F_{\beta^*})$ \citep{efron1979bootstrap}.

\subsection{Benchmarks, metrics and Monte Carlo uncertainty}
\label{sec:benchmarks}

Ten polynomial systems on $\R^2$ and $\R^3$ serve as the main benchmarks (all
defined in the released code). Four have a nonzero symmetry algebra inside the
affine class, of dimension $r^\star$: the Hopf normal form $\dot x = x - y -
x(x^2+y^2)$, $\dot y = x + y - y(x^2+y^2)$ ($r^\star = 1$), the planar rotation
$\dot x = -y$, $\dot y = x$ ($r^\star = 2$), the Euler equations of a symmetric
top with moments of inertia $(1,1,2)$ ($r^\star = 2$) and the radial flow onto
the sphere $\dot x = x(1-|x|^2)$ on $\R^3$ ($r^\star = 3$,
$\mathfrak{so}(3)$). Six have none: Van der Pol ($\dot x = y$, $\dot y = -x + y
- x^2y$), Duffing ($\dot x = y$, $\dot y = x - x^3 - y/5$), Lotka--Volterra
($\dot x = x - xy$, $\dot y = xy - y$), Sel'kov ($\dot x = 1/10 - x + y/2 +
x^2y$, $\dot y = -y/2 - x^2y$), the Euler equations of an asymmetric rigid body
with moments $(1,2,3)$, and Lorenz with parameters $10$, $28$ and $8/3$. Each is
fitted in the polynomial class of its own degree (1 for the rotation, 2 for the
rigid bodies, Lotka--Volterra and Lorenz, 3 otherwise) with affine candidate
generators. The target measure $\nu$ is uniform on the box $[-s,s]^n$, with $s
= 2$ for Duffing and Sel'kov, $s = 20$ for Lorenz and $s = 1$ otherwise
(except in the selected-generator study of \SFtightness e, where all four
systems are on $[-1,1]^2$, and in the trajectory study of \STtrajectory, where
Van der Pol is on $[-3,3]^2$ and Hopf on $[-1.5,1.5]^2$); the states are drawn from $\nu$ and the errors are
Gaussian. The calibration grid
crosses $N \in \{100, 200, 400, 800, 1600, 3200\}$ with noise levels $\{0.01,
0.03, 0.1, 0.3\}$. The broken-Hopf family is $F_\varepsilon = F_{\mathrm{Hopf}}
+ \varepsilon x^2\partial_x$ (the term $\varepsilon x^2$ enters $\dot x$ only);
the rotation defect grows continuously from zero with $\varepsilon$. Every system has rational coefficients, so its exact symmetry
algebra is the exact rational nullspace of the bracket map, computed with
symbolic arithmetic, and the true defect of any generator follows in closed
form from monomial moments; the algebra dimensions are cross-checked against
an independent characterisation (for linear systems, the dimension of the
commutant of the coefficient matrix). For the dictionary families the true
algebra is the numerical nullspace of the exact bracket map, with a gap of
more than ten orders of magnitude between its zero and nonzero singular values.

The truncation study fits Van der Pol, Duffing, Hopf and Sel'kov on their
boxes in the polynomial classes of degree 1, 2 and 3 (the correct class, as a
control), with $N \in \{400, 1600\}$, noise level $0.02$ and $300$
replications per setting; the lack-of-fit test compares each fit with the
cubic class. The sparsity study fits the broken-Hopf members with $\varepsilon
\in \{0, 0.02, 0.05, 0.1, 0.2, 0.4\}$ on $[-1,1]^2$ in the cubic class by
sequentially thresholded least squares (STLSQ): a least-squares fit followed
by ten rounds in which every monomial coefficient of magnitude below the
threshold is set to zero and each component is refitted on the remaining
monomials, with thresholds $0$ (plain least squares), $0.05$, $0.1$ and
$0.25$, $N = 1000$, noise level $0.05$ and $300$ replications per setting; the
threshold rule is applied to the STLSQ fit and \symcert{} to the unpruned
least-squares fit of the same data. The noise study near the tolerance uses
broken-Hopf members with $\varepsilon$ tuned so that the smallest true defect
in the affine class is $1.2$, $1.5$, $2$ or $4$ times $\delta$, and the
hardening Duffing oscillator $\dot x = y$, $\dot y = -x - 0.3x^3 - 0.2y$
(smallest true defect $0.102$), all on $[-1,1]^2$ in the cubic class, with $N
\in \{100, 400, 1600\}$, noise levels $\{0.05, 0.1, 0.2, 0.3\}$ and $300$
replications in each of the $60$ settings; as no generator is within
tolerance, every declaration by the threshold rule at $\delta$ and every
certification is false.

\emph{False certification} is the event that the declared subspace contains a
generator whose true defect exceeds $\delta$; because the declared object is a
subspace, this is evaluated as a generalised eigenvalue of the true forms
restricted to it, not by sampling directions. \emph{Detection} is the event
that the declared dimension is at least the true one, and \emph{recovery} the
event that, in addition, the declared subspace contains the true algebra up to
a largest principal angle of $10^\circ$ in the $L^2(\nu)$ geometry of generator
fields. Unless stated otherwise $\delta = \alpha = 0.05$.

Where an interval is given for a rate estimated from $k$ events in $T$
replications, it is the exact Clopper--Pearson interval
\citep{clopper1934use}; when no event occurs, the one-sided $95\%$ upper bound
is $1 - 0.05^{1/T}$ ($0.010$ for $T = 300$). For a rate given without an
interval, the $95\%$ Clopper--Pearson interval extends at most $0.058$ from the
estimate for $T \ge 300$, $0.072$ for $T \ge 200$ and $0.116$ for $T \ge 80$,
the smallest number of replications used.
Simultaneous statements over $S$ settings use the Bonferroni level $0.05/S$
per setting. These intervals quantify Monte Carlo error only. The validity of
the certificate rests on \Cref{prop:simultaneous}; the simulations check that
the implementation and the assumptions behave as the theorem says, and probe
what happens when the assumptions fail.

\subsection{Ocean drifter data}
\label{sec:driftermethods}

Hourly positions and velocities are from the Global Drifter Program product of
Elipot et al.~\citep{elipot2016global}, version 2.01~\citep{elipot2022hourly}, obtained from
the NOAA AOML ERDDAP server on 26 October 2026 for 30--45$^\circ$N,
180--140$^\circ$W and 1 June to 30 September of 2016--2019. We kept rows
recorded before the loss of the drogue and with interpolation gaps of at most
six hours, split each drifter's record into contiguous hourly segments of at
least $150$ hours, removed a centred running mean ($49$ hours unless stated) to
isolate the near-inertial band, and took central differences for the
acceleration, dividing by the local Coriolis parameter so that time is measured
in units of $1/f$; velocities are expressed in units of their root-mean-square.
Every step acts identically on both velocity components and commutes with
rotations, so the processing introduces no preferred direction and cannot
remove a rotation invariance that the data have. Removing the running mean
does discard the sub-inertial flow (mean, geostrophic and Ekman currents),
which need not be rotation invariant, so the certificate concerns the dynamics
of the high-passed, near-inertial velocities, not of the full flow. The model classes are polynomials of degree 1 to 3 in the velocity
with affine candidate generators, the target measure is the standard Gaussian
in the normalised velocity (extrapolation factor $1.02$ for the linear class),
and the regions are OLS, White, cluster-robust over drifters ($31$ clusters)
and Newey--West within drifters, applied to each drifter's segments
concatenated in time order. The processed subset is distributed with the
code.

\subsection{Reporting checklist}
\label{sec:checklist}

A certificate is interpretable only together with the choices that give the
tolerance its meaning. We recommend reporting, in this order:
\begin{enumerate}\itemsep1pt
\item the target measure $\nu$ (the region the claim is about) and the output
  metric $W$, or the characteristic scales used to nondimensionalise the state;
\item the candidate class of generators;
\item the model class or dictionary, how it was chosen (fixed in advance, by
  \Cref{prop:classes}, or by a heuristic), and any model-error bound with its
  justification;
\item the observation model and the confidence region used (exact, White,
  cluster-robust, Newey--West), with the regime in which it is valid;
\item the extrapolation factor $\sqrt{\kappa\kappa'}$ of the observed states;
\item $U(\theta)$ for every generator reported, the tolerance $\delta$ and
  whether it was fixed before the analysis, the dimension of the certified
  plug-in subspace, and, if it is at least two, its closure diagnostic;
\item the lower bound $\min_\theta L(\theta)$ when the claim is that no
  symmetry exists.
\end{enumerate}

\subsection{Proofs}
\label{app:proofs}

\paragraph{\Cref{lem:props}.} Write $u = DF\,\xi$ and $v = D\xi\,F$ as elements
of $L^2(\nu)$ with the inner product induced by $W$. Then $\defect^2 =
\|u-v\|^2/(\|u\|^2+\|v\|^2)$. The inequality $\|u-v\|^2 \le
2(\|u\|^2+\|v\|^2)$ gives $\defect \le \sqrt 2$; $\defect = 0$ iff $u = v$,
i.e.\ $[\xi,F]=0$ $\nu$-almost everywhere, which for continuous fields means on
the support of $\nu$. When the support is the closure of an open set, the
bracket vanishes on its interior, which is equivalent to the flow of $\xi$
mapping solutions to solutions there \citep[Ch.~9]{lee2013smooth}; for
polynomial or other analytic fields a bracket that vanishes on an open set
vanishes identically. On a lower-dimensional support no such conclusion
follows. Under
$\xi \mapsto c\xi$ both $u$ and $v$ scale by $c$, under $F \mapsto c'F$ by
$c'$, so the ratio is invariant, and hence also under any invertible linear
reparameterisation of the candidate class, because $\defect$ depends on
$\theta$ only through $\xi_\theta$. For $x = Az+b$ the pulled-back fields are
$\tilde F(z) = A^{-1}F(Az+b)$ and $\tilde\xi(z) = A^{-1}\xi(Az+b)$, with
$D\tilde F(z) = A^{-1}DF(Az+b)A$ and likewise for $\xi$, so $\tilde u(z) =
A^{-1}u(Az+b)$ and $\tilde v(z) = A^{-1}v(Az+b)$. With $\tilde\nu$ the
transported measure and $\tilde W = A^{\!\top}WA$, $\int \tilde u^{\!\top}\tilde
W\tilde u\,d\tilde\nu = \int u^{\!\top}Wu\,d\nu$ and likewise for $v$ and $u-v$,
so the ratio is unchanged. With $\tilde W = W = I$ this requires
$A^{\!\top}A \propto I$, i.e.\ a similarity; for any other $A$ the scalings of
the components differ and the ratio changes in general (\STunits). \qed

\paragraph{\Cref{lem:bilinear}.} $DF_\beta$ and $F_\beta$ are linear in $\beta$,
and $\xi_\theta$ and $D\xi_\theta$ are linear in $\theta$, so each of
$DF_\beta\,\xi_\theta$ and $D\xi_\theta F_\beta$ is bilinear in $(\theta,\beta)$
pointwise. Their $\|\cdot\|_{\nu,W}$ norms are therefore quadratic forms in
$\beta$ for fixed $\theta$ and in $\theta$ for fixed $\beta$, and the symmetric
square roots of the Gram matrices give the matrices of the lemma.
$\defect^2 = \theta^{\!\top}C\theta / \theta^{\!\top}D\theta$; if $D\theta = 0$
then $M_\theta\beta = 0$, both terms vanish, $L_\theta\beta = 0$ and the defect
is undefined. On the range of $D$ the minimum of the Rayleigh quotient is the
smallest generalised eigenvalue. \qed

\paragraph{\Cref{prop:simultaneous}.} On the event $E = \{\beta^\star \in
\mathcal B\}$ we have $L(\theta) \le \defect_\nu(\xi_\theta;F_{\beta^\star}) \le
U(\theta)$ for every $\theta$ with a defined defect, by definition of the
infimum and supremum. Since $\Prob(E) \ge 1-\alpha$ and $E$ depends neither on
$\theta$ nor on $\delta$, the statement holds simultaneously. For
\eqref{eq:fwer}: on $E$, any $\theta$ and $\delta$ with $U(\theta) \le \delta$
have $\defect_\nu(\xi_\theta;F_{\beta^\star}) \le \delta$, whatever rule
produced them, so the bad event is contained in $E^c$. \qed

\paragraph{\Cref{prop:closedform}.} (i) With $\beta = \bhat + \hat\Sigma^{1/2}u$
each norm is $\|a + Bu\|$ with $a = L_\theta\bhat$, $B =
L_\theta\hat\Sigma^{1/2}$ (or the same with $M_\theta$), and its extrema over
the ball are trust-region subproblems; dividing the supremum of the numerator
by the infimum of the denominator bounds the supremum of the ratio. (ii) The
denominator being positive on the ball, $\sup_u
\|L_\theta\beta(u)\|^2/\|M_\theta\beta(u)\|^2 \le t$ iff
$t\|M_\theta\beta(u)\|^2 - \|L_\theta\beta(u)\|^2 \ge 0$ for all $u$ in the
ball, iff its minimum over the ball is non-negative; the minimum of a
quadratic over a ball has no duality gap, even when the quadratic is
indefinite \citep{more1983computing,conn2000trust}, so the dual bound decides
the sign up to the bisection tolerance. The lower bound $L(\theta)$ is obtained
in the same way from the sign of the minimum of $\|L_\theta\beta(u)\|^2 -
t\|M_\theta\beta(u)\|^2$. (iii) In model~\eqref{eq:model} with a fixed design
$\Phi$ of full column rank, $\bhat -
\beta^\star \sim N(0, \sigma^2 \Sigma_0)$ with $\Sigma_0 = I_n \otimes
(\Phi^{\!\top}\Phi)^{-1}$ and $\mathrm{RSS}/\sigma^2 \sim
\chi^2_{\mathrm{dof}}$ independent of $\bhat$, so
$\|\hat\Sigma^{-1/2}(\bhat-\beta^\star)\|^2/Q$, with $\hat\Sigma =
\hat\sigma^2\Sigma_0$, is exactly $\mathrm{F}_{Q,\mathrm{dof}}$ distributed. \qed

\paragraph{\Cref{prop:lof}.} The $F$-statistic compares the drop in residual
sum of squares from the fitted class to $\mathcal H$ with the residual mean
square of $\mathcal H$; with Gaussian errors and the truth in $\mathcal H$ its
numerator is a noncentral chi-square with noncentrality equal to the squared
norm, summed over rows and components and divided by $\sigma^2$, of the part
of $F$ in $\mathcal H$ orthogonal to the fitted class over the design, which is
$N\|h\|^2_{\mu_N}/\sigma^2$. The power is increasing in $\lambda$, which gives
the first claim. Because $F_{\beta^\circ}$ is the least-squares projection over
the design, $h$ lies in $\mathcal H_\perp$, so $\|[\xi,h]\|_{\nu,W} \le
G(\xi)\|h\|_{\mu_N}$ by~\eqref{eq:invisible-gain}, and by the triangle
inequality the numerator $\|[\xi,F_{\beta^\circ}+h]\|_{\nu,W}$ differs from
$\|[\xi,F_{\beta^\circ}]\|_{\nu,W}$ by at most this amount. Let $h^\star$
attain the maximum in
\eqref{eq:invisible-gain} and scale it to $\|h^\star\|_{\mu_N} =
\sigma(\lambda_\gamma/N)^{1/2}$; then $\|[\xi,h^\star]\|_{\nu,W} =
G(\xi)\sigma(\lambda_\gamma/N)^{1/2}$, and the truth $F_{\beta^\circ}+h^\star$
has the same projection $F_{\beta^\circ}$ and noncentrality $\lambda_\gamma$.
If $[\xi,F_{\beta^\circ}]=0$, the numerator of the defect of the truth is
exactly $\|[\xi,h^\star]\|_{\nu,W}$, which gives~\eqref{eq:invisible}. \qed

\paragraph{\Cref{prop:classes}.} Work in model~\eqref{eq:model} and fix the
design. In class $k < K$ the
least-squares estimate satisfies $\bhat_k - \beta_k^\circ \sim N(0,
\sigma^2\Sigma_{0,k})$, where $\beta^\circ_k$ are the coefficients of the
least-squares projection of $F$ on $\mathcal F_k$ over the design; the residual
sum of squares is $\sigma^2$ times a noncentral chi-square (the noncentrality
coming from the omitted part), independent of $\bhat_k$ and stochastically
larger than the central one. The $F$-statistic of
\Cref{prop:closedform}(iii) is therefore stochastically smaller than
$\mathrm{F}_{Q_k,\mathrm{dof}_k}$, and $\mathcal B_k$ covers $\beta^\circ_k$ with
probability at least $1-\alpha/K$; $\mathcal B_K$ covers $\beta^\star$ with
probability $1-\alpha/K$. On the intersection of these $K$ events, which has
probability at least $1-\alpha$, $F = F_{\beta^\circ_k} + h_k(\beta^\star)$ with
$\|h_k(\beta^\star)\|_{\nu,W} \le \bar\eta_k$ for every $k$, and the $L^2$
model-error bound in class $k$ gives $\defect_\nu(\xi_\theta;F) \le
U_k(\theta)$. \qed

\paragraph{\Cref{prop:lecam}.} With $L_\xi$, $M_\xi$ and $\Delta^\star$ as in the
statement, set $\beta^t =
\beta^0 + t\Delta^\star$. Since $L_\xi \beta^0 = 0$, $\|L_\xi\beta^t\| =
t\,s_{\max}$, and $\|M_\xi \beta^t\| \le \|M_\xi\beta^0\| +
t\|M_\xi\Delta^\star\|$, so $\defect_\nu(\xi;F_{\beta^t}) \ge t s_{\max}/
(\|M_\xi\beta^0\| + t\|M_\xi\Delta^\star\|)$, which exceeds $\delta$ for every
$t > t_\delta = \delta\|M_\xi\beta^0\|/(s_{\max}-\delta\|M_\xi\Delta^\star\|)$.
The Kullback--Leibler divergence between the two data distributions is
$\tfrac{N t^2}{2\sigma^2}\|F_{\Delta^\star}\|^2_{\mu_N} = Nt^2/(2\sigma^2)$, so
by Pinsker's inequality the total variation distance is at most
$t\sqrt N/(2\sigma)$. For any rule, $\Prob_{\beta^0}(\varphi=1) \le
\Prob_{\beta^{t}}(\varphi=1) + t\sqrt N/(2\sigma) \le \alpha + t\sqrt
N/(2\sigma)$ for every $t > t_\delta$; letting $t \downarrow t_\delta$ gives
\eqref{eq:lecam}, with no remainder term. This is the two-point method
\citep[Ch.~2]{tsybakov2009introduction}. \qed

\paragraph{\Cref{prop:kappa}.} Both $\|q\|^2_{\nu}$ and $\|q\|^2_{\mu}$ are
quadratic forms in the coefficient vector $c$ of $q \in \mathcal R$, equal to
$c^{\!\top}G_\nu c$ and $c^{\!\top}G_\mu c$; their ratio is a generalised
Rayleigh quotient whose maximum is
$\lambda_{\max}(G_\mu^{-1/2}G_\nu G_\mu^{-1/2})$, attained at the corresponding
eigenvector. The defect statement follows by applying the bound
$\|q\|^2_\nu \le \kappa\|q\|^2_\mu$ to the numerator, since the bracket lies in
$\mathcal R$, and the reverse bound $\|q\|^2_\mu \le \kappa'\|q\|^2_\nu$ to
each of the two terms of the denominator, both of which lie in $\mathcal R$.
If $G_\mu$ is singular with null
vector $c \ne 0$ then $\|q_c\|_{\mu} = 0$ while $\|q_c\|_{\nu}>0$, so $\kappa =
\infty$; conversely $\kappa = \infty$ forces $G_\mu$ to be singular. \qed

\subsection{Software, reproducibility and compute}
\label{app:repro}

The implementation (Python, NumPy, SciPy and SymPy) and all experiment and
figure scripts are released together with the raw result files (Code
availability). All experiments are deterministic given the seeds recorded in
the code. Every table of numerical results in the manuscript and the
Supplementary Information is written by a script directly from the result
files, and the numerical
results quoted in the running text are asserted against those files by a
checking script, which fails if any has drifted. A test suite verifies the bracket against finite
differences, the closed-form moments against Monte Carlo integration, the
defect against Monte Carlo evaluation of its defining integrals, the exact
symmetry algebra of linear systems against the dimension of the commutant, the
certified trust-region bounds and the sharp ratio bounds against brute-force
search over the ball, the invariance of the defect under changes of units with
a transported metric, the dictionary construction against the closed-form one,
the coverage of the coefficient region and of the certificate against their
nominal levels, and the closure diagnostic on exact Lie algebras. One complete
pass of the experiments costs about $3.6$ CPU-hours as recorded by the
experiment runner. A large language model (Claude, Anthropic) was used to
assist with writing the code, running and analysing the experiments and
drafting the text; the author takes responsibility for the content.
